\subsection[Методика №2: построение и анализ радиолокационных изображений]{\texorpdfstring{Методика №2: построение и анализ радиолокационных\\изображений}{Методика №2: построение и анализ радиолокационных изображений}}
\label{sec:radar-images-method}

В отличие от методики №1, основанной на прямом спектральном анализе S-параметров, методика №2 реализует традиционный для георадиолокации подход, рассмотренный в параграфе 1.2. В его основе лежит переход из частотной области во временную с последующим формированием радиолокационных изображений --- прежде всего A- и B-сканов, --- которые затем могут анализироваться визуально и сопоставляться с геометрией исследуемой среды.

\subsubsection{Обратное преобразование Фурье и формирование A--B-сканов}

Исходные данные представляют собой значения S-параметра
\begin{equation}
    S(f_k, x_j) = |S(f_k, x_j)| e^{i\varphi(f_k, x_j)},
\end{equation}
полученные в дискретных точках частотной сетки $f_k$ (см. параграф 2.2) и на позициях антенны $x_j$. Для каждого положения антенны выполняется обратное дискретное преобразование Фурье, в результате которого восстанавливается временной отклик среды:
\begin{equation}
    s_{\mathrm{in}}(t_n, x_j)
    = \sum_{k=0}^{N-1} S(f_k, x_j)\exp(i2\pi f_k t_n),
\end{equation}
где $t_n = n\Delta t$, а шаг $\Delta t \simeq 1/B$ определяется частотной полосой $B$ радара (см. формулу~\eqref{eq:range-resolution}). Полученный массив $s_{\mathrm{in}}(t_n, x_j)$ образует набор <<сырых>> A-сканов; последовательность таких сканов по координате $x_j$ формирует B-скан.

Перед визуализацией данные проходят обработку для устранения систематических искажений и повышения отношения сигнал/шум.

\paragraph{Удаление фона.}
Для подавления устойчивых паразитных отражений, не связанных с подповерхностными объектами, применяются два взаимодополняющих метода. В первом случае используется измеренный опорный отклик $s_{\mathrm{oa}}(t)$, соответствующий излучению антенны в свободное пространство (режим вычитания open-air). Из каждого A-скана вычитается этот отклик:
\begin{equation}
    s_{\mathrm{rem}}(t) = s_{\mathrm{in}}(t) - s_{\mathrm{oa}}(t),
\end{equation}
что позволяет исключить вклад систематических паразитных составляющих.

Альтернативный способ --- вычитание усредненного по всем трассам отклика (режим вычитания среднего):
\begin{equation}
    s_{\mathrm{rem}}(t) = s_{\mathrm{in}}(t) - \frac{1}{J}\sum_{j=1}^{J}s_{\mathrm{in}}(t, x_j),
\end{equation}
где $J$ --- число сканов в серии. Метод эффективно устраняет горизонтальные артефакты, вызванные антенным <<звоном>> и стабильными поверхностными отражениями. Однако в случае однородных слабоотражающих слоев применение этой формулы может привести к нежелательному подавлению полезной информации. Поэтому выбор метода зависит от условий измерений и конкретной задачи.

\paragraph{Коррекция нулевого момента времени.}
После выполнения обратного преобразования Фурье временной отсчет $t = 0$ может не совпадать с физически значимой плоскостью антенн. Положение истинного <<нуля>> уточняется с использованием калибровочных измерений, выполненных с металлическим отражателем, расположенным на известных расстояниях. На основе этих данных определяется постоянная задержка $t_0$, по которой сдвигается временная шкала:
\begin{equation}
    s_{t_0}(t) = s_{\mathrm{rem}}(t - t_0).
\end{equation}

\paragraph{Усиление амплитуды.}
С увеличением глубины амплитуда отраженных сигналов снижается вследствие геометрического расхождения волнового фронта и затухания в материале. Для компенсации этого эффекта применяется усиление с использованием множителей $g_j(t)$, позволяя повысить видимость отражений от глубоко расположенных объектов:
\begin{equation}
    s_{\mathrm{gain}}(t) = \sum_j g_j(t)s_{t_0}^{j}(t),
\end{equation}
где $g_j(t)$ --- весовая функция $j$-го временного окна (например, линейная или экспоненциальная).

Последовательное выполнение этих операций дает информативные A-сканы. Их совокупность в координатах $(x, t)$ образует B-скан, пригодный для визуального анализа и интерпретации.

В ряде прикладных сценариев могут быть полезны и дополнительные процедуры, упомянутые в параграфе 1.2, однако в данной работе они рассматривались как вспомогательные.
